Redistricting reform faces a fundamental collective action problem: the legislators who must enact reform are the same officials who benefit most from the status quo. We develop a formal political economy model showing that incumbents treat favorable district lines as a durable asset yielding an ``incumbency rent'' in each election cycle. This rent creates strong incentives to block reform through ordinary legislative channels. We identify three bypass channels through which reform can succeed despite incumbent resistance: direct democracy (ballot initiatives), judicial intervention (state and federal courts striking partisan maps), and federal legislation (Congress exercising its Article I \S4 authority). Applying the model to 35 reform attempts across 22 states between 1990 and 2024, we show that reform probability increases significantly with state-level electoral competitiveness, the availability of the ballot initiative mechanism, and the existence of active gerrymandering litigation. States with all three conditions succeeded in reform at a rate of 78\%, compared to 14\% for states with none. The analysis frames Tracks P.1--P.5 of this research program: each paper examines one bypass channel in depth.
